\clearpage
\section{The precise stabilizer-action contract and normalizer obstruction}\label{app:stabilizer_contract}

Evidence: direct audit of the original statement, a new elementary
reduction, and an explicit actual-source/actual-type H1--H2 witness.
No Lean closure, class-wide arithmetic
realization, or H5 witness is claimed.

\textbf{Primary audit verdict:} incomplete, with the smallest missing
implication being the identification of the declared finite action with
the action intended on a newly constructed common arithmetic candidate.
The original text does not define that identification by a common matrix
normalizer.

\hypertarget{stabilizer_contract-exact-contract-and-what-is-already-an-input}{%
\subsection{1. Exact contract, and what is already an
input}\label{stabilizer_contract-exact-contract-and-what-is-already-an-input}}

The controlling text is
\path{closure_attempt162/inputs/original_reading.txt}:

\begin{itemize}
\tightlist
\item
  Proposition 9.8, starting line 4834: a finite group acting by
  automorphisms of the one-break Clifford datum acts symplectically on
  its intrinsic envelope.
\item
  Definition~11.1, lines 4911--4923: a bridge supplies an equivariant
  isometry to an actual type quotient for a finite cyclic stabilizer C,
  with common central character and additional dyadic lift data.
\item
  Proposition 11.3 and Theorem~11.5: once the two honest cyclic actions
  are chosen, their discrepancy is a character, killed by equality of phases at the generator.
\item
  Theorem~15.4 and Theorem~16.2: breakwise comparisons and transports are
  supplied and tensored. They do not construct a common arithmetic host
  from separate actual type targets.
\end{itemize}

Thus H2 and H3 already settle the finite action comparison \textbf{on
each supplied break target}. They do not leave a second free stabilizer
action to invent on that same target. Their transport into a new common
arithmetic type is another assertion. Neither a map C into its matrix
normalizer nor arbitrary transported operators are identified with that
assertion by the original text.

The cyclic group in Definition~11.1 is not explicitly required there to
have order prime to p.~The reduction below assumes that order condition. It applies to the declared tame prime-to-p case, and cannot silently be
applied to an arbitrary finite-order Frobenius operator.

\hypertarget{stabilizer_contract-full-scalar-extension-rigidity-for-a-native-tame-action}{%
\subsection{2. Full scalar-extension rigidity for a native tame
action}\label{stabilizer_contract-full-scalar-extension-rigidity-for-a-native-tame-action}}

Let an actual maximal simple type in GL\_n(F) have a totally ramified
parameter field K/F of degree n=p\^{}r and residue matrix rank one. Let

\begin{Verbatim}[fontsize=\small,breaklines=true,breakanywhere=true]
W=J¹/H¹,  P=(C××J¹)/{(theta(h)^(-1),h):h in H¹},
\end{Verbatim}

and let eta be its irreducible Heisenberg representation. Let C be
finite of order prime to p.~Suppose the C action on P fixes the scalar
center and is genuinely induced by character-preserving matrix
normalizers of this type. It is enough that it is an honest action whose
action on W comes from those normalizers. Then:

\begin{enumerate}
\def\labelenumi{\arabic{enumi}.}
\tightlist
\item
  C acts trivially on W.
\item
  C acts trivially on the full scalar extension P.
\item
  Every honest C action Omega on eta implementing this action is
  Omega(c)=chi(c)I for one character chi of C.
\end{enumerate}

\textbf{Proof.} The retained normalizer theorem in Block162 gives

\begin{Verbatim}[fontsize=\small,breaklines=true,breakanywhere=true]
N=K×J¹,          N/(F×J¹) = Z/nZ.
\end{Verbatim}

The action on W factors through this cyclic p-group, so a homomorphic
image of the prime-to-p group C is trivial.

For any central extension P of the elementary abelian p-group W by C×,
the kernel of

\begin{Verbatim}[fontsize=\small,breaklines=true,breakanywhere=true]
Aut_(C×)(P) -> Aut(W)
\end{Verbatim}

is naturally Hom(W,C×). Indeed, an automorphism phi fixing the quotient
and the center has phi(x)=a(x\_bar)x. The scalar a is independent of the
lift x and multiplication shows that a is a character. Conversely every
such character defines an automorphism. Since W has exponent p, this
kernel is a finite elementary abelian p-group. The action of C lies in
this kernel and therefore is trivial. Finally an implementing operator
commutes with the irreducible P action, so is scalar by Schur. Honesty
makes its scalar a character. QED.

This strengthens the earlier W-level obstruction: passing to the full
scalar pushout cannot hide a nontrivial native prime-to-p action above
an identity quotient action.

\hypertarget{stabilizer_contract-exact-attachment-criterion-in-this-native-setting}{%
\subsection{3. Exact attachment criterion in this native
setting}\label{stabilizer_contract-exact-attachment-criterion-in-this-native-setting}}

Suppose a marked source scalar extension P\_src has a symplectic
embedding into W, and an actual finite sector of eta realizes that
extension. Fix its center-preserving attachment and a Heisenberg
intertwiner T. If the host has the native tame action of Section~2, the
attachment intertwines the full honest C actions if and only if

\begin{Verbatim}[fontsize=\small,breaklines=true,breakanywhere=true]
the source action on P_src is trivial, and
chi_src = chi_host.
\end{Verbatim}

Necessity follows from equivariance and Section~2. For sufficiency,
trivial extension actions make both honest carrier actions scalar by
Schur, and equal characters commute through every Heisenberg
intertwiner. Any sector is stable under the scalar host action.

If the source C action on its symplectic quotient is trivial, its action
on the full scalar extension is already trivial by the Hom-kernel
argument. For several break factors, the source character is the product
of their characters. Thus the criterion is a single equality of scalar characters, rather than a further system of symplectic placement
equations.

This applies to the original C=\{1\} fixture immediately. For nontrivial
C, H3 equates source and \textbf{supplied break-target} phases. It does
not equate their product with a separately built common host's phase
without a dictionary. We do not change a fixed canonical beta extension
to enforce this equality.

\hypertarget{stabilizer_contract-h2-and-matching-carrier-dimensions-do-not-force-trivial-c}{%
\subsection{4. H2 and matching carrier dimensions do not force trivial
C}\label{stabilizer_contract-h2-and-matching-carrier-dimensions-do-not-force-trivial-c}}

Here is an actual H1--H2 witness under the literal original allowance of
separate type targets. It is not a claimed H5-admitted input, and does
not assert equal ambient matrix ranks for source and break target.

Take F=Fp((t)), p odd. The retained primary source by Tsushima,
Definition~3.1, Lemmas 2.6, 2.8, 3.9, 3.12, 3.14 and Corollary 3.15, with
R(X)=X\^{}p and m=1, constructs a genuine rank-p Weil representation
irreducible on wild inertia. Its wild image is extraspecial of order p³.
The source envelope is a symplectic plane, and its actual tame image
C=mu\_(p+1) acts by multiplication on the roots of X\^{}(p²)+X. This
action is nontrivial and has order p+1. The source's finite Heisenberg
carrier has dimension p.

For the actual type target, take the unramified quadratic extension E/F
and the maximal order A=End\_OF(OE). Let beta=t\^{}(-2)a, where a is a
unit generating the residue extension. This is a minimal simple stratum
of even level two. The standard minimal simple-character formulas give

\begin{Verbatim}[fontsize=\small,breaklines=true,breakanywhere=true]
J¹=U_A¹,  H¹=U_E¹ U_A²,
W_theta=End_Fp(Fp²)/Fp²,     dim_Fp W_theta=2,
dim eta_theta=p.
\end{Verbatim}

These are actual groups of a simple character, not an abstract
replacement. The type residue units act by conjugation. Write
sigma(x)=x\^{}p and decompose

\begin{Verbatim}[fontsize=\small,breaklines=true,breakanywhere=true]
End_Fp(Fp²)=Fp² direct-sum Fp² sigma.
\end{Verbatim}

On the quotient represented by b sigma, conjugation by u in Fp²× is

\begin{Verbatim}[fontsize=\small,breaklines=true,breakanywhere=true]
b -> u^(1-p)b.
\end{Verbatim}

Its image is the norm-one torus, of order p+1. Consequently the actual
projective residue action of Fp²×/Fp× on the two-dimensional quotient is
nontrivial. It has the same symplectic C-module as the source: identify
the cyclic groups using the displayed ratio map. The two invariant
alternating forms differ by a nonzero Fp scalar, and multiplication by
lambda in Fp²× rescales an alternating form by Norm(lambda).
Surjectivity of the norm gives an equivariant symplectic isometry.

This supplies the actual source, actual type quotient, matching degree-p
Heisenberg carriers, and nontrivial equivariant H2 action. For p=3, C
has order 4 and the action is an irreducible rotation of the F3 plane.

The target here is in GL\_2(F), while the source has rank p.~This is
permitted by the literal breakwise definitions, which contain no
relation between the individual type targets' ambient ranks. Block161
already recorded that exact omission. If a stricter intended contract
requires every break target in a specified GL\_(p\^{}r)(F), this example
does not settle that stricter question. We have not supplied H5 for a
common candidate or certified fixed canonical beta-extension phases for
this example. It refutes only the proposed inference that actual H2 plus
equality of finite carrier dimensions forces C trivial.

\hypertarget{stabilizer_contract-result-for-the-next-construction-attempt}{%
\subsection{5. Result for the next construction
attempt}\label{stabilizer_contract-result-for-the-next-construction-attempt}}

There are three distinct cases:

\begin{longtable}[]{@{}
  >{\raggedright\arraybackslash}p{(\columnwidth - 2\tabcolsep) * \real{0.5000}}
  >{\raggedright\arraybackslash}p{(\columnwidth - 2\tabcolsep) * \real{0.5000}}@{}}
\toprule\noalign{}
\begin{minipage}[b]{\linewidth}\raggedright
Given contract
\end{minipage} & \begin{minipage}[b]{\linewidth}\raggedright
Exact remaining issue
\end{minipage} \\
\midrule\noalign{}
\endhead
\bottomrule\noalign{}
\endlastfoot
Original C=\{1\} fixture & No stabilizer issue remains. \\
Native prime-to-p C action in a degree-p\^{}r totally ramified host &
Source extension action must be trivial. Then compare one scalar character. \\
Nontrivial source symplectic C action under the original finite contract
& A dictionary for the arithmetic action on a common host is needed. Requiring an
inner normalizer makes this impossible for the preceding host. Transporting external operators alone was already done in Block 161. \\
\end{longtable}

Thus ``solve arbitrary nontrivial stabilizer attachment by adjusting the
current totally ramified binder'' is not a viable native-inner target.
That is a precise failure of that stronger construction route, not a
refutation of 16.2.

\hypertarget{stabilizer_contract-evidence-and-self-review}{%
\subsection{Evidence and self-review}\label{stabilizer_contract-evidence-and-self-review}}

Dependencies: original Sections~9--11,15--16; Block162's source-backed
normalizer theorem; Block165's base realization contract and
same-candidate attachment; the retained published Tsushima text; the
retained minimal-stratum group formulas. The new Hom-kernel lemma and
attachment criterion are ordinary elementary proofs. The witness uses
the displayed arithmetic constructor inputs, not a numerical experiment.
No independent review of this new note is claimed in this file, and no
new Lean compilation was attempted.

Self-review explicitly checked: the prime-to-p qualification, honest
action versus projective implementers, center-fixed scalar extension
automorphisms, separate target rank versus original rank, and the
absence of an H5 witness.
